\chapter{Conclusion and Contributions}
\label{chap:conclusion}

\section{Study Conclusion}
This thesis developed and evaluated a comparative framework for explainable Bayesian survival prediction. It connected coefficient-prior choice, repeated predictive evaluation, calibration, computational diagnostics and posterior interpretation within a common study. Normal, horseshoe and continuous spike-and-slab priors shared a piecewise-exponential proportional-hazards likelihood and predictor representation. CoxPH, random survival forest and DeepSurv placed their performance in the context of established survival methods.

The analysis covered three clinical cohorts with different endpoints and information structures: GBSG2, WHAS500 and METABRIC. Five repeats of five-fold outer cross-validation yielded \ConfirmatoryPlanned{} original model--fold outcomes, with \ConfirmatoryComplete{} accepted predictions. A further \RobustnessPlanned{} outcomes examined controlled changes in information, follow-up, predictors and model specification. Training-only preprocessing and matched evaluation subjects made the within-cohort comparisons directly interpretable.

The main empirical result is that Normal and continuous spike-and-slab produced closely aligned probability predictions across all three cohorts. Their paired IBS differences ranged from -0.000008 to +0.000061. The complete original horseshoe comparisons in WHAS500 were also small relative to their uncertainty. Larger predictive differences occurred between model families in METABRIC, where RSF and DeepSurv had lower IBS than the Normal-prior model in the secondary paired estimates. The training-information analyses then identified a specific change in relative horseshoe benefit in WHAS500.

Posterior coefficients, hazard ratios, credible intervals and directional probabilities provided a direct interpretation of clinical associations within accepted models. Grouped posterior contributions and exact linear SHAP often produced similar predictor rankings. Their sampling assessment showed that agreement between targets and stability across explanation samples are separate empirical properties. These findings form a coherent account of what the evaluated survival models predict, how their priors behave and how their fitted associations can be interpreted.

\section{Answers to the Research Questions}
\subsection{RQ1: Predictive Behaviour Under Different Coefficient Priors}
The controlled prior comparison found small predictive differences in the clinical predictor sets studied. Normal minus continuous spike-and-slab IBS was +0.000025 in GBSG2, +0.000061 in WHAS500 and -0.000008 in METABRIC. All five complete primary intervals included zero. The original horseshoe arm supplied 55 accepted outcomes across the three cohorts, with its 20 diagnostic outcomes reflected in the coverage of the corresponding comparisons. The estimates place prior choice alongside computational acceptance when assessing the practical behaviour of a Bayesian survival model.

The secondary comparisons add clinical context to this result. Normal and CoxPH had closely aligned IBS in every cohort. In METABRIC, Normal minus RSF IBS was +0.007956, with a corrected 95\% interval of [+0.004531, +0.011381]. The corresponding DeepSurv contrast was +0.002893 [+0.000707, +0.005079]. WHAS500 showed a different descriptive ordering, with CoxPH and the complete Bayesian arms giving lower mean IBS than the forest and neural models. Thus the evidence supports model selection within a defined clinical setting and evaluation design.

\subsection{RQ2: Behaviour Under Reduced Information}
The controlled information changes identified a cohort-specific shrinkage response. At 25\% training, the Normal minus horseshoe IBS contrast increased by +0.006323 in WHAS500, with a descriptive 95\% interval of [+0.002512, +0.010561]. METABRIC gave a smaller change of +0.000211 [-0.001149, +0.001508], while GBSG2 supplied three of ten accepted pairs. Continuous spike-and-slab changes remained close to zero under the same training reductions.

This result gives a precise answer to the information-regime question: the relative effect of shrinkage depends on both the prior and the clinical setting. It also shows the value of varying training information within the same cohort. That design separates the controlled perturbation from the wider differences in case mix, endpoints and predictors between cohorts.

\subsection{RQ3: Posterior Interpretation, Agreement and Sampling Stability}
The Bayesian models supplied uncertainty-aware summaries for named predictor associations. Positive lymph nodes and progesterone receptor were prominent in GBSG2; age and heart rate were prominent in WHAS500; age at diagnosis and Nottingham Prognostic Index were prominent in METABRIC. Coefficients, hazard ratios and directional probabilities expressed the sign, magnitude and uncertainty of these associations on their modelled scales.

The grouped contribution comparison extended this interpretation across priors. Same-model Spearman agreement with linear SHAP ranged from 0.859 to 0.966. Of 1,640 evaluable sampling checks, 1,538 met the correlation and top-five thresholds and 102 did not, so the all-comparisons requirement for H3a and H3b was not met. The resulting distinction is substantive: two explanation summaries can agree positively while some rankings remain sensitive to the subjects, backgrounds and seeds used to estimate them.

\subsection{RQ4: Dependence on Predictor and Model Specification}
The robustness analyses quantified prior contrasts under changes to follow-up, feature definitions, missing-data handling, baseline intervals and prior scales. Normal minus continuous-spike-and-slab differences generally remained small, with condition-specific changes in sign and precision. In METABRIC, the contrast stayed close to zero under Nottingham Prognostic Index components, stage inclusion and five-imputation prediction averaging. The reported pair counts connect each estimate to its contributing evaluations.

Calibration and proportional-hazards checks provided further information about the selected model structures. Calibration-in-the-large was near zero on average for the complete Bayesian arms, while slopes varied across cohorts. Training-fit proportional-hazards checks identified repeated departures for selected METABRIC predictors. These results place coefficient interpretation and predictive performance within the fitted survival structure rather than reducing robustness to a single score.

\section{Research Contributions}
\subsection{A Controlled Framework for Comparing Survival Priors}
The first contribution is an integrated experimental framework in which coefficient priors can be compared under a common survival likelihood, predictor set and validation design. This makes the prior contrast interpretable while retaining the practical context supplied by CoxPH, RSF and DeepSurv. The framework combines probability error, discrimination and calibration with computational acceptance, so that the choice of prior can be assessed through several properties of the resulting model.

\subsection{Empirical Evidence Across Three Clinical Settings}
The second contribution is a repeated-validation benchmark spanning three clinical survival settings. It quantifies the small Normal versus continuous-spike-and-slab differences, the cohort-specific comparator results and the changes observed under controlled information reductions. The additional predictor, censoring and model-specification analyses provide a structured account of how those effects behave when the analysis conditions change.

\subsection{A Common Basis for Uncertainty-Aware Interpretation}
The third contribution connects posterior effect interpretation with a common grouped contribution target across priors. Coefficients, hazard ratios and directional probabilities describe fitted associations; grouped contributions describe the magnitude of predictor participation in the modelled risk function. Exact linear SHAP supplies a mathematically explicit comparison target. Evaluating target agreement separately from sampling stability makes the meaning of each explanation result clear.

\subsection{Computational Sensitivity as Part of Model Assessment}
The post-hoc sampler analysis adds a distinct computational result. Increasing target acceptance to 0.999 achieved diagnostic acceptance for all 75 horseshoe evaluations, with all nine resulting prior contrasts estimable. Their intervals included zero and their Holm-adjusted $p$ values were one. This analysis shows how sampling configuration can change computational reliability while leaving the measured predictive differences small. Its separate reporting preserves the interpretation of both computational settings.

\section{Implications for Survival-Modelling Practice}
The findings support evaluating a Bayesian survival model through the combination of its predictions, probability calibration, posterior summaries and sampling behaviour. Models with similar IBS can offer different parameter uncertainty and computational properties. Models with similar discrimination can also differ in the level or spread of predicted probabilities. These properties answer related but distinct questions and can be presented together within a reproducible study.

Prior choice is most informative when its effect is evaluated under shared clinical information and validation conditions. Here, the complete Normal and continuous spike-and-slab models produced very similar predictions. The horseshoe information-regime result in WHAS500 identifies a more specific setting in which relative shrinkage behaviour changed. This gives a concrete basis for considering the information available to estimate the model, rather than choosing a prior solely by its general reputation for sparsity or flexibility.

For interpretation, posterior uncertainty and feature contribution serve complementary purposes. A credible interval describes uncertainty about a model coefficient conditional on the fitted data and assumptions. A contribution ranking describes relative participation in a specified prediction function. Sampling checks then show how consistently that ranking is reproduced when its estimation sample changes. Keeping these targets explicit supports clear communication of Bayesian survival results.

\section{Final Synthesis}
Across these cohorts, coefficient priors often changed posterior summaries and sampling behaviour more than average probability error. The Bayesian analysis brings these properties together through predictive performance, calibration and uncertainty about fitted associations. The controlled comparisons and explanation analysis provide a concrete basis for choosing and communicating Bayesian survival models for a defined prediction task.
